\chapter{Literature Review}
\section{Image-based plant disease classification}
Hughes and Salathé released the PlantVillage collection of labelled plant-leaf photos to support research on automatic diagnosis~\cite{hughes2015dataset}. Mohanty et al. trained deep networks on it and reached very high accuracy, while noting that photos taken in real fields are harder~\cite{mohanty2016plant}. Ferentinos tested several CNN designs on a larger multi-crop collection~\cite{ferentinos2018plant}. Too et al. compared fine-tuned networks and found DenseNet performed best~\cite{too2019plant}. Each of these studies uses its own classes and data splits, so their numbers are not directly comparable with ours.

Kamilaris and Prenafeta-Boldú review the wider use of deep learning in agriculture~\cite{kamilaris2018survey}. PlantDoc collects leaf photos taken outside the laboratory~\cite{singh2020plantdoc}, and Gui et al. study how background and lighting in field photos affect recognition~\cite{gui2021field}. These works show that field photos are an important next step after laboratory datasets.

\section{Transfer learning and efficient networks}
EfficientNet grows a network in a balanced way by scaling its depth, width, and input image size together~\cite{tan2019efficientnet}. We use its smallest version, B0, and tune how it is trained. MobileNetV2 and MobileNetV3 are other compact networks designed for phones~\cite{sandler2018mobilenet,howard2019mobilenet}.

Yosinski et al. show that early layers of a trained network learn general features, while later layers become specific to the original task~\cite{yosinski2014transfer}. This means the number of layers we retrain (the fine-tuning depth) is an important choice. Keeping the whole pretrained network fixed is a simple, common baseline. Retraining more layers can adapt better to leaves, but costs more time. We let the search decide.

\section{Hyperparameter and evolutionary search}
Grid search tries every combination from a fixed list of values; random search simply tries randomly chosen settings. Bayesian optimisation builds a model of which settings are likely to work well~\cite{snoek2012bayesian}. Hyperband saves time by stopping weak trials early~\cite{li2016hyperband}. Random search is a strong and simple competitor, so we use it as our control: it tells us whether DE's guided search adds value over picking the same number of settings at random.

Evolutionary methods have been used successfully to design and tune networks. Real et al. used an evolutionary method to design image classifiers~\cite{real2019evolution}. Wang et al. used particle swarm optimisation and a hybrid DE method to design CNNs~\cite{wang2018psoarchitecture,wang2018dearchitecture}. DEHB combines DE with Hyperband for tuning~\cite{awad2021dehb}. These results motivate applying DE to our seven training settings.

\begin{table}[htbp]\centering\small
\caption{Published results, each measured on its own dataset and protocol.}
\begin{tabularx}{\textwidth}{@{}>{\raggedright\arraybackslash}p{.19\textwidth}>{\raggedright\arraybackslash}X >{\raggedright\arraybackslash}p{.23\textwidth}r@{}}\toprule
Study & Dataset / model & Search or comparison & Accuracy\\\midrule
Mohanty et al.~\cite{mohanty2016plant} & PlantVillage, 38 classes; AlexNet/GoogLeNet & Training/model variants & 99.35\%\\
Ferentinos~\cite{ferentinos2018plant} & 58 categories; CNNs & CNN architectures & 99.53\%\\
Too et al.~\cite{too2019plant} & PlantVillage, 38 classes; DenseNet & Fine-tuned architectures & 99.75\%\\
Gui et al.~\cite{gui2021field} & Field-PV; adapted CNN & Background/features & 72.03\%\\\bottomrule
\end{tabularx}
\end{table}

\section{Research gap}
Most earlier plant-disease studies choose training settings by hand. Few apply an evolutionary search to these settings and compare it fairly with a fixed baseline and with random search at the same cost. This project fills that gap for tomato leaves: it uses fair data splits, a baseline fixed in advance, the same number of trials for each search method, and repeated final training runs. It also links the algorithm to a working offline demo.
